% GENERATED FILE. Edit canonical lessons, then run npm run generate:books.
% canonical-source-hash: fnv1a64:0f4939bf74a4fd8d
% canonical-lessons: TA-C166-vilakkerru, TA-C166-veti, TA-C166-poruttukkol, TA-C166-camatanam-cey, TA-C166-yoci

\chapter{To light a lamp, To burst, To put up with, To make peace, To consider}
\label{ch:ta-to-light-a-lamp-to-burst-to-put-up-with-to-make-peace-to-consider}
\hlchaptermodality{166}

\begin{chapteropening}
\textbf{By the end of this chapter:} \emph{I can say to light a lamp, to burst, to put up with, to make peace and to consider.}

Complete the last lesson of chapter 166: I can say to light a lamp, to burst, to put up with, to make peace and to consider.
\end{chapteropening}

\section[\texorpdfstring{\ta{விளக்கேற்று} (viḷakkēṟṟu)}{விளக்கேற்று (viḷakkēṟṟu)}]{\ta{விளக்கேற்று} (viḷakkēṟṟu) — to light a lamp}

\label{lesson:TA-C166-vilakkerru}



\begin{quote}
Before the new one: say the Tamil for to dance, then the Tamil for to worship.
\end{quote}

\subsection*{You'll want to know: \ta{விளக்கேற்று}}
\textbf{\ta{விளக்கேற்று}} — \emph{viḷakkēṟṟu} — \textquotedblleft{}to light a lamp\textquotedblright{}.

\begin{cousinweb}
\textbf{\ta{விளக்கேற்று}} (\emph{viḷakkēṟṟu}) is native: \emph{viḷakku}, \textquotedblleft{}lamp\textquotedblright{} (DEDR 5496), and \emph{ēṟṟu}, \textquotedblleft{}to raise, to light\textquotedblright{}, the causative of \emph{ēṟu}, \textquotedblleft{}to rise\textquotedblright{} (DEDR 916).
\end{cousinweb}

\begin{grammarlens}[title={What you've built}]
One more word for this chapter.
\end{grammarlens}

\subsection*{Your turn}
\begin{itemize}
\raggedright
  \item \emph{Say it:} \emph{viḷakkēṟṟu}
  \item \emph{Say it:} \emph{viḷakkēṟṟu}, once more
  \item \emph{Say it:} say \emph{viḷakkēṟṟu}
  \item \emph{Recall:} say the Tamil for to dance, then the Tamil for to worship, then say \emph{viḷakkēṟṟu} again
\end{itemize}

\begin{tcolorbox}[breakable,colback=teal!4,colframe=teal!35!black,title={Before you move on}]
What does \ta{விளக்கேற்று} mean? (\textquotedblleft{}To light a lamp\textquotedblright{}.) Say it once more.
\end{tcolorbox}
\section[\texorpdfstring{\ta{வெடி} (veṭi)}{வெடி (veṭi)}]{\ta{வெடி} (veṭi) — to burst, to explode}

\label{lesson:TA-C166-veti}



\begin{quote}
Before the new one: say the Tamil for to worship, then the Tamil for to light a lamp.
\end{quote}

\subsection*{You'll want to know: \ta{வெடி}}
\textbf{\ta{வெடி}} — \emph{veṭi} — \textquotedblleft{}to burst, to explode\textquotedblright{}.

\begin{cousinweb}
\textbf{\ta{வெடி}} (\emph{veṭi}) is native Dravidian. DEDR 5473 lists Malayalam \emph{veṭi}, \textquotedblleft{}explosion\textquotedblright{}.
\end{cousinweb}

\begin{grammarlens}[title={What you've built}]
One more word for this chapter.
\end{grammarlens}

\subsection*{Your turn}
\begin{itemize}
\raggedright
  \item \emph{Say it:} \emph{veṭi}
  \item \emph{Say it:} \emph{veṭi}, once more
  \item \emph{Say it:} say \emph{veṭi}
  \item \emph{Recall:} say the Tamil for to worship, then the Tamil for to light a lamp, then say \emph{veṭi} again
\end{itemize}

\begin{tcolorbox}[breakable,colback=teal!4,colframe=teal!35!black,title={Before you move on}]
What does \ta{வெடி} mean? (\textquotedblleft{}To burst, to explode\textquotedblright{}.) Say it once more.
\end{tcolorbox}
\section[\texorpdfstring{\ta{பொறுத்துக்கொள்} (poṟuttukkoḷ)}{பொறுத்துக்கொள் (poṟuttukkoḷ)}]{\ta{பொறுத்துக்கொள்} (poṟuttukkoḷ) — to put up with, to tolerate}

\label{lesson:TA-C166-poruttukkol}



\begin{quote}
Before the new one: say the Tamil for to light a lamp, then the Tamil for to burst.
\end{quote}

\subsection*{You'll want to know: \ta{பொறுத்துக்கொள்}}
\textbf{\ta{பொறுத்துக்கொள்}} — \emph{poṟuttukkoḷ} — \textquotedblleft{}to put up with, to tolerate\textquotedblright{}.

\begin{cousinweb}
\textbf{\ta{பொறுத்துக்கொள்}} (\emph{poṟuttukkoḷ}) is native: \emph{poṟuttu}, \textquotedblleft{}having borne\textquotedblright{}, from \emph{poṟu}, \textquotedblleft{}to bear, to tolerate\textquotedblright{} (DEDR 4565; Malayalam \emph{poṟukka}), and \emph{koḷ}, \textquotedblleft{}to take\textquotedblright{} (DEDR 2151).
\end{cousinweb}

\begin{grammarlens}[title={What you've built}]
One more word for this chapter.
\end{grammarlens}

\subsection*{Your turn}
\begin{itemize}
\raggedright
  \item \emph{Say it:} \emph{poṟuttukkoḷ}
  \item \emph{Say it:} \emph{poṟuttukkoḷ}, once more
  \item \emph{Say it:} say \emph{poṟuttukkoḷ}
  \item \emph{Recall:} say the Tamil for to light a lamp, then the Tamil for to burst, then say \emph{poṟuttukkoḷ} again
\end{itemize}

\begin{tcolorbox}[breakable,colback=teal!4,colframe=teal!35!black,title={Before you move on}]
What does \ta{பொறுத்துக்கொள்} mean? (\textquotedblleft{}To put up with, to tolerate\textquotedblright{}.) Say it once more.
\end{tcolorbox}
\section[\texorpdfstring{\ta{சமாதானம்} \ta{செய்} (camātāṉam cey)}{சமாதானம் செய் (camātāṉam cey)}]{\ta{சமாதானம்} \ta{செய்} (camātāṉam cey) — to make peace, to calm someone down}

\label{lesson:TA-C166-camatanam-cey}



\begin{quote}
Before the new one: say the Tamil for to burst, then the Tamil for to put up with.
\end{quote}

\subsection*{You'll want to know: \ta{சமாதானம்} \ta{செய்}}
\textbf{\ta{சமாதானம்} \ta{செய்}} — \emph{camātāṉam cey} — \textquotedblleft{}to make peace, to calm someone down\textquotedblright{}.

\begin{cousinweb}
\textbf{\ta{சமாதானம்} \ta{செய்}} (\emph{camātāṉam cey}) is \emph{camātāṉam}, \textquotedblleft{}peace, reconciliation\textquotedblright{}, from Sanskrit \emph{samādhāna}, \textquotedblleft{}putting together, reconciliation\textquotedblright{} (CDIAL 13210), with native \emph{sey}, \textquotedblleft{}to do\textquotedblright{} (DEDR 1957).
\end{cousinweb}

\begin{grammarlens}[title={What you've built}]
One more word for this chapter.
\end{grammarlens}

\subsection*{Your turn}
\begin{itemize}
\raggedright
  \item \emph{Say it:} \emph{camātāṉam cey}
  \item \emph{Say it:} \emph{camātāṉam cey}, once more
  \item \emph{Say it:} say \emph{camātāṉam cey}
  \item \emph{Recall:} say the Tamil for to burst, then the Tamil for to put up with, then say \emph{camātāṉam cey} again
\end{itemize}

\begin{tcolorbox}[breakable,colback=teal!4,colframe=teal!35!black,title={Before you move on}]
What does \ta{சமாதானம்} \ta{செய்} mean? (\textquotedblleft{}To make peace, to calm someone down\textquotedblright{}.) Say it once more.
\end{tcolorbox}
\section[\texorpdfstring{\ta{யோசி} (yōci)}{யோசி (yōci)}]{\ta{யோசி} (yōci) — to consider, to think something over}

\label{lesson:TA-C166-yoci}



\begin{quote}
Before the new one: say the Tamil for to put up with, then the Tamil for to make peace.
\end{quote}

\subsection*{You'll want to know: \ta{யோசி}}
\textbf{\ta{யோசி}} — \emph{yōci} — \textquotedblleft{}to consider, to think something over\textquotedblright{}.

\begin{cousinweb}
\textbf{\ta{யோசி}} (\emph{yōci}) is a verb made from \emph{yōcaṉai}, \textquotedblleft{}thought, consideration\textquotedblright{}, a loan from Sanskrit \emph{yojanā}, \textquotedblleft{}a joining, an arranging\textquotedblright{}.
\end{cousinweb}

\begin{grammarlens}[title={What you've built}]
One more word for this chapter.
\end{grammarlens}

\subsection*{Your turn}
\begin{itemize}
\raggedright
  \item \emph{Say it:} \emph{yōci}
  \item \emph{Say it:} \emph{yōci}, once more
  \item \emph{Say it:} say \emph{yōci}
  \item \emph{Recall:} say the Tamil for to put up with, then the Tamil for to make peace, then say \emph{yōci} again
\end{itemize}

\begin{tcolorbox}[breakable,colback=teal!4,colframe=teal!35!black,title={Before you move on}]
What does \ta{யோசி} mean? (\textquotedblleft{}To consider, to think something over\textquotedblright{}.) Say it once more.
\end{tcolorbox}
